\section*{Chapitre \ref{ch:m3:natural-gas-and-lng} --- Gaz naturel et GNL}

\begin{solution}{exo:m3:natural-gas-and-lng:1}
$10/0.293071/1.16 = \qty{29.42}{EUR/MWh}$. Un therm vaut 0.1 MMBtu, donc 1~\$ le therm, soit
$100/1.30 = 76.9$ pence par therm.
\end{solution}

\begin{solution}{exo:m3:natural-gas-and-lng:2}
La demande hivernale dépasse la production et doit être satisfaite à partir du stockage rempli en été, si bien
le gaz d'hiver doit rémunérer le stockage du gaz d'été. L'arbitrage des propriétaires de stockage le borne
par le haut, au coût du stockage (injection, soutirage, capacité et financement) lorsque
la capacité est libre~; lorsque le stockage est plein ou rare, il peut dépasser ce coût.
\end{solution}

\begin{solution}{exo:m3:natural-gas-and-lng:3}
Le prix variable vaut $1.15 \times 2.00 = \qty{2.30}{\$/\mmbtu}$, supérieur au \omterm{def:m3:physical-commodity-markets:netback}{netback} de
2.10~: il s'annule. Il perd de toute façon les \qty{2.50}{} de frais~; l'enlèvement
perdrait en outre \qty{0.20}{\$/\mmbtu}.
\end{solution}

\begin{solution}{exo:m3:natural-gas-and-lng:4}
Au-dessus de $11.00 + 0.33 = \qty{11.33}{\$/\mmbtu}$.
\end{solution}

\begin{solution}{exo:m3:natural-gas-and-lng:5}
Les détenteurs de capacité de transport par gazoduc au départ du hub~: le spread au-dessus du coût de transport,
\qty{1.70}{\$/\mmbtu}, représente la valeur de leur capacité. Avec cette capacité, on achète au hub, on expédie,
puis on vend à \omterm{def:m3:natural-gas-and-lng:hubs}{Henry Hub} (ou l'on vend la base à terme).
\end{solution}

\begin{solution}{exo:m3:natural-gas-and-lng:6}
Le spread de seuil de rentabilité passe de \qty{0.33}{} à \qty{1.35}{\$/\mmbtu}, soit une hausse de \qty{1.02}{}~: la
traversée plus longue coûte davantage en affrètement et en boil-off.
\end{solution}

\begin{solution}{exo:m3:natural-gas-and-lng:7}
L'Asie l'emporte dans 93 des 125 mois, de mars 2016 à juillet 2026~; le cargo n'est pas
enlevé de février à août 2020 et en février 2021.
\end{solution}

\begin{solution}{exo:m3:natural-gas-and-lng:8}
Le revenu du cargo est le prix de destination diminué du transport et de la regazéification, et son
coût s'élève à 115~\% de \omterm{def:m3:natural-gas-and-lng:hubs}{Henry Hub}, auquel s'ajoutent les frais fixes versés pour la capacité~: la marge est le
\omterm{def:m3:physical-commodity-markets:netback}{netback} diminué des deux. L'essentiel du spread a aussi été capté par celui qui détenait les
droits de destination du cargo, souvent au titre de contrats à long terme signés des années plus tôt, et non par chaque
vendeur~; et les taux de fret ont augmenté avec le spread.
\end{solution}

\begin{solution}{pb:m3:natural-gas-and-lng:1}
\textbf{1.} \qty{3.45}{\$/\mmbtu}. \textbf{2.} \qty{0.54}{\$/\mmbtu}. \textbf{3.}
\qty{0.97}{\$/\mmbtu}. \textbf{4.} 1.4~\% vers l'Europe, 2.5~\% vers l'Asie. \textbf{5.} Le navire
doit revenir à vide (sur lest) avant de pouvoir charger à nouveau.
\textbf{6.} \qty{0.33}{\$/\mmbtu}. \textbf{7.} Asie~: \omterm{def:m3:physical-commodity-markets:netback}{netback} \qty{11.13}{} contre
\qty{10.96}{}. \textbf{8.} $11.13 \times 3.4125$ millions de MMBtu livrés, moins $3.45 \times
3.5$ millions~: environ USD~25.9 millions. \textbf{9.} Le seuil de rentabilité monte à \qty{0.72}{}, au-dessus
du spread de \qty{0.50}{}~: l'Europe. \textbf{10.} Europe~: \omterm{def:m3:physical-commodity-markets:netback}{netback} \qty{68.85}{} contre
\qty{52.62}{\$/\mmbtu}.
\textbf{11.} Les prix aux deux destinations s'écartent au hasard~; le droit de choisir le
plus élevé, net du seuil de rentabilité, est une option sur le spread, qui a une valeur dès que le
spread est volatil. \textbf{12.} Ne pas enlever~: le meilleur \omterm{def:m3:physical-commodity-markets:netback}{netback}, \qty{0.72}{\$/\mmbtu}, était
inférieur au prix variable de \qty{1.87}{}. \textbf{13.} Rien~: les frais sont dus,
que le cargo soit enlevé ou non, donc ils ne font pas partie de la décision marginale. \textbf{14.} C'est
l'acheteur qui en profite (il économise le prix variable diminué du \omterm{def:m3:physical-commodity-markets:netback}{netback})~; le propriétaire de l'installation
de liquéfaction perçoit toujours les frais. \textbf{15.} Fret~: les taux d'affrètement montent, ce qui réduit l'arbitrage.
\textbf{16.} Ses hubs ont payé davantage que l'Asie après la perte du gaz russe par gazoduc, et le
trajet depuis le Golfe américain est plus court. \textbf{17.} Les cargaisons asiatiques en provenance de l'Atlantique doivent suivre
la route la plus longue, ce qui relève le seuil de rentabilité et donc la prime que l'Asie doit payer.
\textbf{18.} Pour la stabilité des prix liée à ses autres coûts, la sécurité d'approvisionnement, ou parce que
les contrats indexés sur le pétrole étaient ce que le vendeur proposait. \textbf{19.} \emph{Résultat clé~:} l'Asie
doit payer \qty{0.33}{\$/\mmbtu} de plus que l'Europe (\qty{1.35}{} si la traversée dure 50 jours).
\textbf{20.} Les cargaisons de GNL flexibles, qui se dirigent vers le meilleur \omterm{def:m3:physical-commodity-markets:netback}{netback} jusqu'à ce que les spreads tombent au
coût du transport maritime.
\end{solution}

\begin{solution}{iq:m3:natural-gas-and-lng:1}
\omterm{def:m3:natural-gas-and-lng:hubs}{Henry Hub} est une interconnexion physique de gazoducs en Louisiane et le point de livraison
des contrats à terme américains~; le TTF est un \omterm{def:m3:natural-gas-and-lng:hubs}{point d'échange virtuel} du réseau néerlandais où tout gaz du
système peut changer de mains. Les unités et les devises diffèrent aussi.
\iqlookfor{hub physique ou virtuel, et les unités.}
\end{solution}

\begin{solution}{iq:m3:natural-gas-and-lng:2}
La demande de gaz est dominée par le chauffage et la production d'électricité, deux usages saisonniers, tandis que la production est
constante et le stockage coûteux~; la demande de pétrole est moins saisonnière, et le pétrole se stocke à bon marché dans des réservoirs.
\iqlookfor{demande saisonnière face à une production constante et au coût du stockage.}
\end{solution}

\begin{solution}{iq:m3:natural-gas-and-lng:3}
Chaque mois, une option d'achat de gaz à 115~\% de \omterm{def:m3:natural-gas-and-lng:hubs}{Henry Hub} et de revente au meilleur
hub de destination diminué du transport et de la regazéification~: une option sur spread (sur le maximum de
plusieurs spreads), dont la prime est constituée des frais fixes.
\iqlookfor{option sur spread, prix variable comme strike, frais comme prime.}
\end{solution}

\begin{solution}{iq:m3:natural-gas-and-lng:4}
Un élargissement brutal de la base~: une panne de gazoduc ou un pic de demande locale éloigne le hub local de
\omterm{def:m3:natural-gas-and-lng:hubs}{Henry Hub}, au-delà du coût de transport~; les événements météo (un gel) font bouger les deux jambes de
montants différents~; appels de marge sur les jambes papier.
\iqlookfor{base et capacité, et non le niveau des prix.}
\end{solution}

\begin{solution}{iq:m3:natural-gas-and-lng:5}
Ils avaient vendu leur production à terme sur les bourses~; lorsque les prix ont été multipliés par dix, les
contrats à terme vendeurs exigeaient une marge de variation en espèces, alors que le gaz physique qui les compensait
ne serait vendu que plus tard. La couverture était saine économiquement, mais elle consommait une liquidité
dont ils ne disposaient pas.
\iqlookfor{liquidité contre solvabilité, marge sur les couvertures cotées en bourse.}
\end{solution}

\begin{solution}{iq:m3:natural-gas-and-lng:6}
Données d'entrée~: positions et calendriers des cargaisons, contraintes contractuelles, courbes à terme par hub,
courbes de fret et état des canaux. Pour chaque cargaison, calculer le \omterm{def:m3:physical-commodity-markets:netback}{netback} vers chaque destination
faisable et à chaque date, choisir sous réserve des créneaux de terminal et de la disponibilité des navires (un
problème d'affectation), et présenter le choix, la marge et sa sensibilité aux spreads.
\iqlookfor{\omterm{def:m3:physical-commodity-markets:netback}{netbacks}, contraintes, optimisation d'affectation.}
\end{solution}
